\section{RESULTS}
\label{sec:results}

The question admits two decomposed sub-questions. First: can a research method of applying transformative game design theory to this problem domain be demonstrated? Second: does that research method yield expected results? Each is addressed in turn.

\subsection{Demonstrating the Method: The Design Process as Contribution}

The first sub-question is answered by the thesis itself. The systematic progression from theoretical analysis through iterative design to empirical testing constitutes a worked example of operationalizing transformative game design theory in debate. The reasoning connecting each design iteration to observed player behavior is documented in Annexes A, B, and G; whether or not the resulting format succeeds transformatively, the method has been rendered explicit and repeatable.

\subsection{Evaluating the Results: Evidence from Playtests}

The second sub-question concerns whether the designed format shows signs of achieving its intended effects. Two playtests were conducted and analyzed using the clue-taxonomy method described in Section~\ref{sec:methods}; the complete taxonomy is in Annex~I. Each of the five evaluation fields corresponds to a prediction made in Section~\ref{sec:design} on the basis of the theoretical framework; the following analysis evaluates how those predictions fared under empirical testing. Consistent with the clue taxonomy's weighting principles, behavioral and transcript-confirmed evidence is treated as stronger than self-report throughout; the meaningful absence of predicted negative clues is also treated as evidentiary, particularly where it holds under conditions of emotional intensity.

\subsubsection{Field 1: Engagement, Appeal, and Long-Term Viability}

For a format to achieve any transformative aims, it must first generate genuine voluntary engagement: without it, the transformational container cannot form, and the repeated participation that enables bleed and skill development cannot occur. The Inner Voices framing and removal of win/loss outcomes were accordingly predicted to generate this engagement by replacing status-competition with imaginative investment; Section~\ref{sec:design} also identified format novelty as a potential barrier to entry, predicting that onboarding cost would be real but not prohibitive.

The strongest positive evidence for this prediction is behavioral. In the second playtest, participants spontaneously adopted the term ``inner voices'' (a term absent from the ruleset) to refer to each other during play. Organic language generation of this kind is coded as a strong clue in Annex~I: it suggests genuine imaginative investment that outran the format's explicit scaffolding. Supporting behavioral evidence includes animated post-session conversations extending beyond the structured debrief in both sessions, and players experimenting with creative choices (accents, performance styles, personal props) that were neither instructed nor rewarded by any formal mechanism. Self-reported engagement is consistent with this picture: participants described the format as ``surprisingly liberating'' and called the Huddle phase ``very very fun.'' These reports constitute moderate evidence given the playtest context and social desirability concerns specific to Field~1 (see Annex~I preamble), but they align with the behavioral pattern.

Two negative engagement clues were observed. The scripted Rolling Out ritual, reported as ``very robotic'' in the first playtest, recurred as a disengagement signal in the second, with participants mumbling scripted phrases and requesting justification. This is a strong negative clue (multi-session, behavioral, immersion-breaking) and is addressed separately as a design gap in Field~3, where its mechanistic implications are most directly felt. The commendation system produced a subtler negative: evaluators in both sessions found commendation assignment difficult and cognitively overwhelming at small informal scale. This is a calibration issue specific to context: the system was designed for tournament use where Evaluators would have prior familiarity with the badge list, but it represents a genuine negative engagement signal for this playtest context.

The differential appeal across participant profiles deserves direct framing rather than treatment as an ambiguity. Three participants across both sessions (all competitive men with strong traditional-format track records) expressed reservations, one describing the format as ``anticlimactic'' and missing the ``ego-boosting'' experience of decisive victory. Women and debaters who had been less successful in traditional formats expressed consistent preference. This pattern confirms the design's intent: the format was not designed to serve the highly competitive profile that traditional formats already serve well, and the differential response is evidence of the mechanic operating as designed rather than as a limitation. The 40-minute onboarding brief required in both sessions confirms the learning curve is real; both sessions produced substantively engaged play within a single session, consistent with the prediction that cost is one of onboarding rather than fundamental inaccessibility.

The prediction was confirmed at behavioral level: genuine imaginative investment manifested in emergent behavior the ruleset did not script, while the differential engagement pattern is consistent with design intent. The commendation and Rolling Out issues are design gaps, not evidence against the core prediction.

Engagement is a necessary but insufficient condition for the format's transformative aims: whether the removal of win/loss outcomes also removed the productive intellectual challenge that drives growth is what Field~2 addresses.

\subsubsection{Field 2: Productive Pressure Mechanisms}

The transformative potential of debate operates through challenge: the cognitive and social demands of the activity are what push participants to grow, and a format that reduces anxiety by also reducing intellectual rigor would substitute comfort for development. Section~\ref{sec:design} explicitly identified whether competitive motivation could survive without win/loss outcomes as one of the format's central design tensions, and the prediction was that productive pressure would persist, redirected rather than removed, driven by the commendation architecture replacing victory as the evaluative anchor.

The evidence confirms that pressure did persist, at high levels. In the first session, players constructed elaborate and nuanced solutions, made careful value-priority arguments, and delivered notably passionate speeches. In the second, players spontaneously produced contextually rich, roleplay-oriented speeches with worldbuilding detail, none of which was instructed. Participants in the first session described ``architecting solutions'' as ``very complex and interesting strategically,'' a self-report consistent with the behavioral evidence. Both sessions also produced strikingly high levels of acting and roleplay; the risk this raised (whether emotional expressiveness came at the cost of analytical rigor) did not materialize, as elaborate solutions, practicalities-based reasoning, and careful value hierarchies remained present alongside the emotional engagement. In the second playtest, participants self-identified polite counter-argumentation as a format-native growth edge, initiating an unprompted group conversation about how to diminish another Inner Voice's contribution without being dismissive: a growth discussion that standard debate contexts do not produce because the skill is not rewarded there.

The prediction held in outcome but not in mechanism. The mechanic driving effort across both sessions was not primarily the commendation system but a social one: participants wanted to be recognized and appreciated by their fellow Inner Voices and by the Subject for the quality and resonance of their contributions. The commendation architecture provided the evaluative frame, but the operative motivator was emergent social recognition, a more nuanced and arguably healthier form of competitive pressure than binary outcome-seeking, tied intrinsically to the quality of engagement rather than to a win. This is an informative deviation from the design prediction: it suggests the commendation system may be instrumentally secondary to the social dynamics it enables, and that future design iterations should attend to those dynamics rather than to the formal evaluative structure alone.

The Subject role produced a persistent design gap across both iterations. In the first playtest, the Subject player described not knowing ``what to say,'' uncertain how to hold emotional investment and evaluative neutrality simultaneously. This gap in S0 and S1 instructions was addressed in the second iteration; but a second Subject scaffolding issue emerged in the second playtest: the Subject ran out of prepared Huddle questions before the allotted time elapsed, revealing that the role requires more preparation scaffolding than the ruleset provided. The second session also offered a partial resolution to the design tension about Subject cognitive burden: the Subject who navigated the role most successfully was the one who leaned into roleplay rather than maintaining analytical detachment. Simply inhabiting the character, responding emotionally, and asking from genuine curiosity resolved the cognitive tension between roleplay and evaluative functions simultaneously. This is consistent with the broader argument in Section~\ref{sec:background} that debate is fundamentally a form of role-playing. The absence of apathy across both sessions (no participant reduced effort in the absence of win/loss outcomes) is a meaningful negative clue confirming that the prediction's outcome holds.

Productive pressure was confirmed as preserved; the mechanism by which it operates is social rather than evaluative (an informative deviation from design prediction), and the Subject scaffolding gap represents an identified design challenge, partially addressed, that persists across iterations.

With pressure established as present and redirected via social dynamics, Field~3 asks whether it is redirected toward useful growth: toward skills and perspectives that transfer beyond the debate context.

\subsubsection{Field 3: Growth Direction -- Useful Skills and Perspectives}

Sustained challenge is a precondition for growth; growth direction determines what is actually being developed. A format that engages and pressures without developing transferable skills fails the transformative aim at its center. The format's procedural rhetoric (specifically the Inner Voices roles, the S0 bias disclosure, and the Subject framing) was accordingly predicted to elicit perspective-taking, contextual reasoning, and collaborative discourse without explicit instruction, because the rules themselves structurally reward these behaviors. The prediction, derived from \textcite{bogost2007persuasive}'s account of procedural rhetoric, is that no participant should need to be told to take other perspectives; the format was designed to make doing otherwise unnatural.

The strongest evidence comes from the second playtest and takes the form of three distinct behavioral clusters, each constituting an independent confirmation of the prediction. First: all players immediately adopted second-person address, direct eye contact with the Subject, and contextual worldbuilding detail from the outset of the debate, none of which was instructed in the ruleset. This spontaneous perspective-taking orientation is behavioral and coded as a strong clue in Annex~I. Second: players organically developed a polite and collaborative form of counter-argumentation, shifting focus and building on others' contributions rather than dismantling them; participants noted that ``outright attacking others' points feels wrong, disrespectful, and unnecessary.'' A norm that arose from the format's structural logic, not from any rule. Third: across the second round of speeches, players moved spontaneously toward alignment on a common solution, converging rather than intensifying conflict: a behavioral inversion of standard debate's production pattern, and the clearest instance of procedural rhetoric functioning as designed. These three clusters are independent: each corresponds to a different aspect of the prediction (perspective-taking, collaborative framing, synthesis orientation) and each emerged without instruction.

Supporting evidence from the first playtest is moderate. Participants articulated learning to ``think about the context of decisions,'' to be ``more human and more honest about human nature,'' and to recognize that ``sometimes decisions are personal, not universal.'' These self-reports align with design intent but are retrospective and unverified by pre-session baseline. A behaviorally notable pattern across both sessions is that play diverged from standard debate conventions as each session progressed: early in the debate participants applied familiar genre habits, but speech by speech they adjusted toward the format's distinct demands. Both sessions also confirmed two negative clues as absent: no non-standard or unhelpful play behaviors were rewarded, and no hollow performative emotionality was observed; the emotional engagement remained substantively connected to reasoning.

The primary negative clue for this field is the Rolling Out ritual. The annex identifies it as the most significant growth-direction design gap: emancipatory bleed requires a sustained transformational container through the closing moment, and a ritual that breaks immersion at the end directly undermines the bleed arc that is the format's central growth mechanism. The ritual was already noted as immersion-breaking in Field~1; its implications are most acute here, where sustained container integrity is a prerequisite for the bleed to complete. The redesigned version in the third iteration (Annex H) remains untested.

The prediction was substantially confirmed: three independent behavioral clusters in the second playtest demonstrate perspective-taking, collaborative framing, and synthesis orientation emerging without instruction, consistent with the procedural rhetoric account. The persistent Rolling Out gap represents an identified design challenge that directly threatens the bleed arc the format depends on and remains the primary open growth-direction issue.

Growth in useful directions is a design-level outcome; Field~4 asks whether the same structural features that produce it also shape the community dynamics, specifically whether the format is a space people want to return to and sustain.

\subsubsection{Field 4: Community Dynamics and Reduced Toxicity}

Even a format that develops useful skills fails if it cannot sustain a community willing to engage with it repeatedly: attrition driven by hostile or hierarchical dynamics limits transformative reach regardless of intrinsic design quality. Removing win/loss and redesigning the evaluative structure toward commendation was accordingly predicted to reduce the adversarial dynamics that contribute to attrition and harm in competitive debate communities. The theoretical framework predicted this as a consequence of restructuring competitive incentives away from zero-sum outcome-seeking: when status derives from demonstrated growth rather than competitive victory, the adversarial dynamics that erode the transformational container are structurally dismantled.

Reduced anxiety is the most consistently confirmed finding across both sessions. Participants reported the format as ``much much less anxiety-inducing'' than standard competitive debate, described a ``feeling of salvation regardless of result,'' and noted that ``everyone distinguished himself by his own style.'' This pattern is very strong (Annex~I): multiple independent reports, consistent across two separate groups, constituting the most overdetermined positive clue in the dataset. One honest caveat applies: the format's lack of established status means less social capital is staked on performance, so some portion of the anxiety reduction may be a novelty artifact rather than a structural feature that would persist in a mature Inside Out circuit. This is an empirical question for future study.

Brave, unconventional play was acknowledged and positively received in both sessions. In the first playtest, one team chose a deliberately vulnerable, emotionally intimate approach (highly uncommon in traditional competitive formats) which was acknowledged in the oral critique. In the second, brave play took multiple forms: one player opened their subject speech with a background soundtrack; the C team built their entire case around a worldbuilt character with emotional stakes; players throughout adopted expressively roleplay-heavy delivery, and the room responded with visible gesture-based approval (laughing, nodding, leaning forward) throughout. The formal commendation system was not consistently invoked for brave play due to calibration immaturity at small scale, but informal social reward was immediate and strong.

The strongest behavioral clue in this field is the emergent norm against adversarial rebuttal. In the second playtest, participants spontaneously articulated that ``outright attacking others' points feels wrong, disrespectful, and unnecessary.'' This norm was not stated in the ruleset; it arose directly from the format's structural logic, a procedural rhetoric outcome in the precise sense \textcite{bogost2007persuasive} theorizes. The mechanism is explicit: no scoring category rewards rebuttal, meaning attacking an argument carries no competitive benefit while building on it does, inverting the rhetorical incentive structure of standard debate. The finding suggests a self-reinforcing alignment between competitive incentives and transformative design intent.

Equally informative is what was not observed. None of the three negative clues for this field (expressions of distress, arrogant or unempathetic behavior, or praise of such behavior) manifested across either playtest or in the transcripts. The absence of distress holds even given the emotional intensity of some moments in the second session, including an unplanned bleed event where a participant shared privately that the scenario resonated with a real-life situation they were facing that same day; the debrief structure provided sufficient containment. The absence of arrogant or dismissive behavior is meaningful given that participants were experienced competitive debaters with established stylistic tendencies in traditional formats.

The prediction was strongly confirmed across two independent sessions: anxiety was consistently reduced, brave unconventional play was rewarded, and the collaborative norm emerged as an unsolicited behavioral outcome of the format's structural logic rather than its rules. The absence of all three negative clues (distress, arrogant behavior, and praise of such behavior) holds across both sessions including moments of genuine emotional intensity, strengthening confidence beyond what positive clues alone would provide.

Community dynamics concern the format as a social space; Field~5 addresses the deepest theoretical prediction: that the format's structural logic would produce epistemic humility as a bleed effect, a measurable shift in how participants hold their own beliefs.

\subsubsection{Field 5: Epistemic Humility}

Epistemic humility is the deepest transformation criterion. The format's  theoretical prediction was that the Inner Voices framing, the S0 bias disclosure, the absence of a winner, and the multi-perspective structure would produce this shift not as a taught skill but as an emergent consequence of inhabiting the format's structural logic.

Moderate supporting evidence appears in the first playtest. One participant explicitly articulated, without prompting: ``every team says valuable stuff and it is important more to see what is important to us.'' This captures precisely the targeted attitudinal shift: the recognition that multiple value frameworks each have genuine validity and that the relevant question is personal priority rather than objective correctness. Additional first-session reflections include the observation that arguing someone is ``morally wrong'' for a personal decision feels ``a bit absurd'' in this format, and that the format makes personal value priorities legible rather than evaluating them against an external standard.

The strongest evidence is behavioral and comes from the second playtest. Synthesis orientation emerged across both rounds of Inner voices speeches, appearing in the first round and intensifying in the second, suggesting the format's structural logic strengthens as the session progresses rather than requiring a settled period before taking effect. This spontaneous drift toward convergence, despite no rule incentivizing it, is transcript-confirmed and constitutes the clearest behavioral indicator in the dataset that the format's procedural logic orients participants toward integration rather than dominance. Players in the second session also organically renamed themselves ``Inner Voices'' during play, a label that positions them as parts of a shared internal deliberation rather than opposing sides; this renaming was subsequently formalized in the third design iteration.

The second playtest debrief generated a particularly theoretically significant conversation about the S0 speech as a structural mechanism. Participants independently identified that the S0 institutionalizes the practice of surfacing personal biases openly before argumentation begins, and contrasted this with standard formats in which the pretense of objectivity is maintained throughout. Those formats, they argued, falsely presuppose that bias can be excluded, propagating the notion that only one type of ``objective'' person can judge well and legitimizing social pressure to perform that identity. The S0 mechanism, by contrast, frames bias as perspective rather than contamination: a structural move toward epistemic humility enforced by the format before any argument is made. That participants arrived at this analysis independently is itself an indicator of the mechanism functioning: the format made its own epistemic logic visible to its players.

The unplanned bleed event (a participant sharing privately that the scenario resonated with a real-life situation they were facing that same day) is the single highest-strength clue across the entire dataset: spontaneous, private (not performed for the researcher), and constituting the transformational mechanism functioning exactly as theorized in Section~\ref{sec:theory}.

No negative clues for this field (unequivocal condemnation of positions or belief-holders, or praise of such condemnation) were observed across either playtest or the transcripts. The absence is meaningful given that both motions involved choices with genuine ethical weight and participants who are practiced in adversarial argumentation.

The prediction was confirmed at proof-of-concept level. The spontaneous synthesis behavior, the ``Inner Voices'' renaming, the S0 debrief conversation, and the unplanned bleed event together constitute behavioral, emergent, and transcript-confirmed evidence that the format produces epistemic humility as the theory predicted: not as instruction, but as an emergent consequence of its structural logic.

\subsection{Proof of Concept Assessment}

Taken together, the five evaluation fields show that the theoretical framework generates testable predictions about debate format behavior, and that those predictions held across two independent playtest sessions. Fields~1, 3, 4, and 5 confirmed their predictions, with the behavioral evidence particularly strong for Fields~4 and 5. Field~2 confirmed the prediction in outcome while revealing an informative deviation in mechanism: social recognition rather than commendations as the operative pressure driver. No field produced evidence contradicting its prediction; the consistent absence of predicted negative clues across both sessions carries meaningful confirmatory weight.

No single piece of evidence is conclusive. Small sample sizes, the absence of longitudinal follow-up, and the researcher's dual role as designer and observer all limit what can be inferred. However, the convergence of multiple independent clues across two separate playtest groups, spanning behavioral, self-report, and transcript-confirmed evidence with predicted negative clues consistently absent, provides a coherent picture that justifies continued investigation.

\subsection{Limitations}

The original design of this study envisioned approximately four iterative playtest sessions. Two were completed within the thesis timeline. While this reduces the number of design iterations and the breadth of participant exposure, both sessions yielded substantive observational and debrief data sufficient for proof-of-concept assessment. Two recurring design gaps (Subject role scaffolding and the Rolling Out ritual's immersion-breaking effect) were partially addressed across iterations but remain open questions; both should be prioritized in the next design cycle. Finally, both playtest groups consisted of experienced debaters already familiar with competitive formats. The format's performance with novice participants (for whom the contrast with standard debate would be less salient) remains untested.
